% !TEX program = pdflatex
% akizumi-paper in a plain article, as under a venue style: commenters and
% final, \anon, abbreviations, table marks, refstyle=short with hyperref.
\documentclass{article}
\usepackage{hyperref}
\usepackage{akizumi-math}
\usepackage[refstyle=short, anonymous]{akizumi-paper}
\newcommenter{jun}{Jun}{teal}
\newcommenter{lance}{Lance}{orange}
\begin{document}

\section{Method}\label{sec:method}
As in \anon{Lin \etal (2025)}, we follow Sutton \etal.
Models, \eg transformers, \ie sequence models, are trained \wrt the loss,
\cf \cref{sec:method}, \cref{fig:x}, \cref{eq:loss}, \cref{eq:loss,eq:two}.
A \jun{is this bound tight?} sentence with a \lance*{rewritten clause} in it.
\begin{equation}\label{eq:loss} L = \E\norm{x}^2 \end{equation}
\begin{equation}\label{eq:two} M = 0 \end{equation}

{\akizumipapersetup{final, anonymous=false}%
Final: a \jun{gone} b, and \lance*{kept plain}; \anon{Lin et al.} named.\par}

\begin{table}[h]\centering
\begin{tabular}{lccc}
\toprule
Method & \makecell{Atari\\100k} & DMC & Sparse \\
\midrule
\multirow{2}{*}{Ours} & \best{0.91} & \second{0.88} & \cmark \\
                      & \second{0.80} & \best{$0.93$} & \xmark \\
\bottomrule
\end{tabular}
\caption{Marks.}
\end{table}
\begin{figure}[h]\centering\rule{2cm}{1cm}\caption{A box.}\label{fig:x}\end{figure}
\end{document}
